\chapter{مقدمه}

% =========================================================
% 1-1
% =========================================================
\section{اهداف و دامنه پروژه}

پیشرفت سریع روش‌های هوش مصنوعی در حوزه‌های پردازش گفتار و
پردازش زبان طبیعی، امکان توسعه سامانه‌هایی را فراهم کرده است که
تعامل انسان و رایانه را از رابط‌های متنی متداول به سمت تعامل طبیعی‌تر
مبتنی بر گفتار هدایت می‌کنند.
در یک سامانه گفت‌وگوی صوتی، کاربر می‌تواند درخواست خود را مستقیماً
از طریق گفتار بیان کند و پاسخ سامانه را نیز به‌صورت صوتی دریافت نماید.
پیاده‌سازی چنین سامانه‌ای نیازمند همکاری چند مؤلفه مستقل در زمینه
پردازش صوت، تشخیص گفتار، درک و تولید زبان و تولید گفتار است.

هدف اصلی این پروژه، طراحی و پیاده‌سازی یک
\textbf{چت‌بات صوتی فارسی}
است که بتواند چرخه کامل دریافت گفتار کاربر، تبدیل آن به متن،
تولید پاسخ متنی و تبدیل پاسخ به گفتار را به‌صورت یکپارچه اجرا کند.
معماری پایه سامانه از سه مؤلفه اصلی تشکیل شده است:
سامانه
\textbf{تشخیص خودکار گفتار}
\LTRfootnote{Automatic \lr{Speech} Recognition (\lr{ASR})}
برای تبدیل صوت ورودی به متن،
یک
\textbf{مدل زبانی بزرگ}
\LTRfootnote{Large Language \lr{Model} (\lr{LLM})}
برای پردازش درخواست و تولید پاسخ،
و سامانه
\textbf{تبدیل متن به گفتار}
\LTRfootnote{\lr{Text}-to-\lr{Speech} (\lr{TTS})}
برای تبدیل پاسخ متنی تولیدشده به گفتار فارسی.

یکی از اهداف اصلی پروژه، امکان اجرای این چرخه به‌صورت محلی
بر روی یک رایانه شخصی است.
اجرای محلی سامانه موجب کاهش وابستگی به سرویس‌های پردازشی خارجی
شده و امکان کنترل بیشتر بر نحوه پردازش داده‌ها، مصرف منابع محاسباتی
و تأخیر سامانه را فراهم می‌کند.
در نتیجه، انتخاب مدل‌ها در این پروژه صرفاً بر اساس کیفیت خروجی
انجام نمی‌شود و معیارهایی نظیر دقت، سرعت استنتاج، مصرف حافظه
و امکان اجرا بر روی سخت‌افزار شخصی نیز مورد توجه قرار می‌گیرند.

برای انتخاب مؤلفه مناسب در هر مرحله، چندین مدل به‌صورت مستقل
مورد ارزیابی قرار گرفته‌اند.
در بخش تشخیص خودکار گفتار، مدل‌هایی بر پایه
\lr{wav2vec 2.0/XLSR}، \lr{Whisper}، \lr{Vosk}
و \lr{NeMo FastConformer}
مورد آزمایش قرار گرفته‌اند.
معماری \lr{wav2vec 2.0} بر یادگیری خودنظارتی بازنمایی‌های گفتاری
استوار است \cite{baevski2020wav2vec2} و نسخه چندزبانه
\lr{XLSR} این رویکرد را به یادگیری بازنمایی میان‌زبانی گسترش
می‌دهد \cite{conneau2020xlsr}.
مدل \lr{Whisper} نیز بر اساس آموزش گسترده چندزبانه و چندوظیفه‌ای
توسعه یافته است \cite{radford2022whisper}.
همچنین معماری \lr{Conformer} با ترکیب سازوکار توجه و کانولوشن،
یکی از معماری‌های مهم مورد استفاده در سامانه‌های جدید تشخیص گفتار
محسوب می‌شود \cite{gulati2020conformer}.

در بخش مدل زبانی، چند مدل قابل اجرای محلی شامل
\lr{Gemma 2}, \lr{Qwen2.5} و \lr{Dorna2/Llama 3.1}
از نظر کیفیت پاسخ، زمان تولید، نرخ تولید توکن و مصرف حافظه
مورد مقایسه قرار گرفته‌اند.
شرط اولیه انتخاب این مدل‌ها آن بود که نسخه چهاربیتی آن‌ها بتواند
به‌طور کامل در حافظه یک پردازنده گرافیکی محلی بارگذاری و بدون
وابستگی به سرویس ابری اجرا شود.
در بخش تبدیل متن به گفتار نیز چند سامانه فارسی از جمله
\lr{Kamtera VITS}، \lr{MMS/VITS} و \lr{Piper}
از نظر کیفیت ادراکی، قابل‌فهم‌بودن گفتار و سرعت تولید صوت
ارزیابی شده‌اند.
روش دقیق ارزیابی و نتایج عددی این مقایسه‌ها در فصل آزمایش‌ها
و نتایج ارائه خواهد شد.

با توجه به ماهیت زنجیره‌ای این معماری، کیفیت سامانه نهایی صرفاً
تابعی از عملکرد مستقل هر مؤلفه نیست.
خطای ایجادشده در یک مرحله می‌تواند به مراحل بعدی منتقل شود.
برای مثال، تشخیص نادرست یک یا چند واژه در مؤلفه \lr{ASR}
ممکن است مفهوم جمله ورودی را تغییر دهد و در نتیجه بر پاسخ
مدل زبانی اثر بگذارد.
به همین ترتیب، حتی در صورت تولید پاسخ متنی مناسب، کیفیت پایین
خروجی \lr{TTS} می‌تواند فهم پاسخ توسط کاربر را دشوار سازد.
ازاین‌رو، علاوه بر ارزیابی مستقل اجزا، ارزیابی عملکرد
انتهابه‌انتهای سامانه نیز یکی از اهداف اصلی پروژه محسوب می‌شود.

یکی دیگر از مسائل مهم در یک سامانه گفت‌وگوی صوتی،
\textbf{تأخیر پاسخ‌گویی}
است.
فاصله زمانی زیاد میان پایان گفتار کاربر و شروع پاسخ صوتی می‌تواند
طبیعی‌بودن مکالمه را کاهش دهد.
بنابراین، در کنار معیارهای کیفیت، معیارهایی مانند
\textbf{ضریب زمان واقعی}
\LTRfootnote{Real-Time Factor (\lr{RTF})}
و زمان تأخیر استنتاج برای مقایسه مدل‌ها مورد استفاده قرار گرفته‌اند.
در مراحل تکمیلی پروژه، تأخیر کل خط لوله شامل
\lr{ASR}، مدل زبانی و \lr{TTS}
نیز به‌صورت انتهابه‌انتها اندازه‌گیری خواهد شد.

معماری سامانه به‌صورت ماژولار در نظر گرفته شده است؛
به این معنا که مؤلفه‌های اصلی تا حد امکان وابستگی مستقیم کمی
به پیاده‌سازی داخلی یکدیگر داشته باشند.
این طراحی امکان جایگزینی یک مدل تشخیص گفتار، مدل زبانی
یا سامانه تبدیل متن به گفتار را بدون بازنویسی کامل خط لوله فراهم
می‌کند.
چنین ساختاری علاوه بر تسهیل آزمایش مدل‌های مختلف، امکان توسعه
و نگهداری ساده‌تر سامانه در آینده را نیز فراهم می‌سازد.

علاوه بر سه مؤلفه اصلی، یک سامانه گفت‌وگوی صوتی کاربردی
به قابلیت‌های تکمیلی دیگری نیز نیاز دارد.
از جمله این موارد می‌توان به
\textbf{تشخیص فعالیت صوتی}
\LTRfootnote{Voice Activity Detection (VAD)}
برای تعیین شروع و پایان گفتار،
پردازش نویز محیط،
پس‌پردازش متن و علائم نگارشی،
و حفظ تاریخچه مکالمه برای پشتیبانی از گفت‌وگوی چندمرحله‌ای اشاره کرد.
این قابلیت‌ها به‌عنوان بخش‌های تکمیلی سامانه در ادامه پروژه
پیاده‌سازی و ارزیابی خواهند شد.

برای تحقق هدف اصلی پروژه، اهداف اجرایی زیر تعریف شده‌اند:

\begin{itemize}

    \item طراحی و پیاده‌سازی یک خط لوله انتهابه‌انتها برای
    تبدیل گفتار فارسی کاربر به پاسخ صوتی فارسی.

    \item ارزیابی و مقایسه چند مدل تشخیص گفتار فارسی
    از نظر دقت، سرعت استنتاج و مصرف منابع محاسباتی.

    \item ارزیابی چند مدل زبانی قابل اجرای محلی
    از نظر کیفیت پاسخ، سرعت تولید و مصرف حافظه.

    \item ارزیابی چند سامانه تبدیل متن به گفتار فارسی
    از نظر کیفیت ادراکی، قابل‌فهم‌بودن گفتار و سرعت تولید.

    \item انتخاب ترکیب مناسب مدل‌های
    \lr{ASR}، \lr{LLM} و \lr{TTS}
    با در نظر گرفتن هم‌زمان کیفیت و هزینه محاسباتی.

    \item طراحی معماری ماژولار به‌منظور فراهم‌کردن امکان
    جایگزینی مستقل مؤلفه‌های اصلی سامانه.

    \item توسعه قابلیت حفظ بافت و تاریخچه مکالمه
    برای پشتیبانی از گفت‌وگوهای چندمرحله‌ای.

    \item بررسی و افزودن پردازش‌هایی مانند تشخیص فعالیت صوتی،
    مدیریت نویز و پس‌پردازش متن برای افزایش پایداری سامانه.

    \item اندازه‌گیری تأخیر هر یک از مؤلفه‌ها و تأخیر
    انتهابه‌انتهای سامانه.

    \item فراهم‌کردن امکان اجرای سامانه بر روی یک رایانه شخصی
    بدون وابستگی دائمی به سرویس‌های پردازشی خارجی.

\end{itemize}

دامنه این پروژه بر تعامل صوتی به زبان فارسی متمرکز است.
هدف پروژه، آموزش یک مدل پایه تشخیص گفتار، مدل زبانی بزرگ
یا مدل تبدیل متن به گفتار از ابتدا نیست؛ بلکه تمرکز اصلی بر
ارزیابی، انتخاب و یکپارچه‌سازی مدل‌های موجود و ایجاد یک سامانه
عملیاتی بر پایه آن‌ها قرار دارد.

همچنین هدف این پروژه توسعه یک دستیار صوتی تجاری در مقیاس بزرگ
یا سامانه‌ای برای پاسخ‌گویی هم‌زمان به تعداد زیادی کاربر نیست.
موضوعاتی مانند زیرساخت ابری گسترده، احراز هویت گوینده،
مقیاس‌پذیری سازمانی و طراحی یک رابط گرافیکی کامل خارج از
دامنه اصلی پروژه قرار دارند.
تمرکز اصلی بر توسعه یک نمونه عملیاتی محلی، ارزیابی علمی اجزای
آن و بررسی قابلیت استفاده از ترکیب مدل‌های انتخاب‌شده در یک
سامانه مکالمه صوتی فارسی است.


% =========================================================
% 1-2
% =========================================================
\section{کاربردها و اهمیت عملی}

تعامل صوتی یکی از طبیعی‌ترین شیوه‌های ارتباط انسان است و استفاده
از آن به‌عنوان رابط میان کاربر و سامانه‌های هوشمند می‌تواند نیاز
به تعامل مستقیم با صفحه‌کلید یا رابط‌های گرافیکی را کاهش دهد.
این ویژگی به‌ویژه در شرایطی که کاربر امکان تایپ‌کردن ندارد یا
استفاده از دست برای تعامل با سامانه دشوار است، اهمیت بیشتری پیدا
می‌کند.

یکی از کاربردهای مستقیم سامانه پیشنهادی،
\textbf{راهنماهای صوتی هوشمند}
است.
در چنین کاربردی، کاربر می‌تواند پرسش خود را به زبان فارسی و
به‌صورت طبیعی بیان کرده و بدون نیاز به مشاهده نمایشگر، پاسخ
صوتی دریافت کند.
این قابلیت می‌تواند در سامانه‌های راهنما، کیوسک‌های اطلاع‌رسانی،
راهنمای خدمات و دستیارهای شخصی مورد استفاده قرار گیرد.

حوزه دیگر،
\textbf{پشتیبانی مشتری}
است.
بخش قابل توجهی از پرسش‌های کاربران در سامانه‌های پشتیبانی شامل
درخواست‌های تکرارشونده و قابل پاسخ‌گویی از طریق یک سامانه
خودکار است.
ترکیب تشخیص گفتار، مدل زبانی و تولید گفتار می‌تواند امکان ایجاد
یک رابط مکالمه‌ای را فراهم کند که پرسش کاربر را دریافت کرده،
پاسخ متنی مناسب را تولید کند و نتیجه را به‌شکل صوتی ارائه دهد.
در صورت اتصال چنین سامانه‌ای به منابع اطلاعاتی اختصاصی،
دامنه کاربرد آن می‌تواند به دستیارهای پشتیبانی تخصصی نیز گسترش یابد.

کاربرد دیگر سامانه در
\textbf{آموزش}
است.
یک رابط مکالمه‌ای فارسی می‌تواند به‌عنوان دستیار آموزشی مورد
استفاده قرار گیرد تا فراگیر بتواند پرسش‌های خود را به‌صورت صوتی
مطرح کند و پاسخ متناسب دریافت نماید.
ماهیت گفت‌وگومحور چنین سامانه‌ای همچنین امکان توسعه تعاملات
چندمرحله‌ای، پرسش و پاسخ و توضیح تدریجی مفاهیم را فراهم می‌کند.

اهمیت عملی این پروژه تنها به کیفیت خروجی مدل‌ها محدود نمی‌شود.
یکی از محورهای اصلی، بررسی امکان اجرای یک سامانه کامل
\lr{ASR--\lr{LLM}--\lr{TTS}}
بر روی سخت‌افزار شخصی است.
مدل‌هایی که کیفیت بالایی دارند ممکن است از نظر حافظه یا زمان
استنتاج برای اجرای محلی مناسب نباشند.
بنابراین، بررسی هم‌زمان کیفیت، سرعت و مصرف منابع می‌تواند به
انتخاب معماری واقع‌بینانه‌تری برای کاربردهای عملی منجر شود.

اجرای محلی همچنین می‌تواند در کاربردهایی که اتصال پایدار به
اینترنت در دسترس نیست یا استفاده دائمی از سرویس‌های ابری مطلوب
نیست، مفید باشد.
ازاین‌رو، دستیابی به یک ترکیب مناسب از مدل‌های نسبتاً سبک،
دقیق و سریع یکی از جنبه‌های کاربردی مهم این پروژه محسوب می‌شود.


% =========================================================
% 1-3
% =========================================================
\section{پرسش‌های پژوهش}

با توجه به اهداف پروژه و ماهیت چندمرحله‌ای سامانه،
پرسش‌های پژوهش به‌گونه‌ای تعریف شده‌اند که هم عملکرد مستقل
مؤلفه‌ها و هم عملکرد سامانه یکپارچه را مورد بررسی قرار دهند.

\subsection*{انتخاب مدل تشخیص گفتار}

\begin{itemize}
    \item
    \textbf{پرسش پژوهش ۱ (\lr{RQ1}):}
    کدام مدل تشخیص خودکار گفتار، مناسب‌ترین توازن میان
    دقت تشخیص، سرعت استنتاج و مصرف منابع محاسباتی را
    برای یک چت‌بات صوتی فارسی محلی فراهم می‌کند؟
\end{itemize}

برای پاسخ به این پرسش، مدل‌های مختلف تشخیص گفتار بر اساس
معیارهایی نظیر
\textbf{نرخ خطای واژه}
\LTRfootnote{Word Error Rate (\lr{WER})}،
\textbf{نرخ خطای نویسه}
\LTRfootnote{Character Error Rate (\lr{CER})}،
ضریب زمان واقعی، تأخیر متوسط و مصرف حافظه گرافیکی
مورد مقایسه قرار می‌گیرند.


\subsection*{انتخاب مدل زبانی}

\begin{itemize}
    \item
    \textbf{پرسش پژوهش ۲ (\lr{RQ2}):}
    کدام مدل زبانی قابل اجرای محلی، مناسب‌ترین توازن میان
    کیفیت معنایی پاسخ، زمان پاسخ‌گویی و مصرف حافظه را
    برای استفاده در سامانه فراهم می‌کند؟
\end{itemize}

در این بخش، مدل‌ها از نظر معیارهای کیفیت متنی و همچنین
ویژگی‌های اجرایی مانند زمان پاسخ، نرخ تولید توکن و مصرف
حافظه مقایسه می‌شوند.


\subsection*{انتخاب مدل تبدیل متن به گفتار}

\begin{itemize}
    \item
    \textbf{پرسش پژوهش ۳ (\lr{RQ3}):}
    کدام سامانه تبدیل متن به گفتار فارسی، مناسب‌ترین توازن
    میان کیفیت ادراکی صوت، قابل‌فهم‌بودن گفتار و سرعت تولید
    را ارائه می‌دهد؟
\end{itemize}

برای پاسخ به این پرسش، کیفیت ادراکی مدل‌های تبدیل متن به گفتار
به همراه معیارهای مرتبط با قابل‌فهم‌بودن خروجی و زمان اجرای
مدل بررسی می‌شود.


\subsection*{عملکرد انتهابه‌انتهای سامانه}

\begin{itemize}
    \item
    \textbf{پرسش پژوهش ۴ (\lr{RQ4}):}
    ترکیب مؤلفه‌های منتخب در یک خط لوله
    \lr{ASR--\lr{LLM}--\lr{TTS}}
    چه تأخیر انتهابه‌انتهایی ایجاد می‌کند و آیا این تأخیر
    برای یک تعامل مکالمه‌ای قابل قبول است؟
\end{itemize}

این پرسش علاوه بر زمان پردازش مستقل هر مؤلفه،
کل فاصله زمانی میان دریافت گفتار کاربر و تولید پاسخ صوتی
سامانه را مورد توجه قرار می‌دهد.


\subsection*{پایداری تعامل و بافت مکالمه}

\begin{itemize}
    \item
    \textbf{پرسش پژوهش ۵ (\lr{RQ5}):}
    افزودن مؤلفه‌هایی مانند تشخیص فعالیت صوتی، پردازش نویز
    و مدیریت بافت مکالمه چه تأثیری بر پایداری و کیفیت تعامل
    با سامانه دارد؟
\end{itemize}

پاسخ نهایی به این پرسش پس از تکمیل مؤلفه‌های مربوطه
و اجرای آزمایش‌های انتهابه‌انتها ارائه خواهد شد.


% =========================================================
% 1-4
% =========================================================
\section{ساختار گزارش}

این گزارش در هشت فصل تنظیم شده است.

\textbf{فصل اول} به معرفی مسئله، اهداف و دامنه پروژه،
کاربردها و اهمیت عملی، پرسش‌های پژوهش و ساختار کلی گزارش
اختصاص دارد.

\textbf{فصل دوم} مفاهیم نظری مورد نیاز برای درک سامانه
را معرفی می‌کند.
در این فصل مبانی پردازش گفتار، تشخیص خودکار گفتار،
مدل‌های زبانی بزرگ، معماری \lr{Transformer}،
تبدیل متن به گفتار، تشخیص فعالیت صوتی، مدیریت بافت مکالمه
و معیارهای ارزیابی مرتبط بررسی می‌شوند.

\textbf{فصل سوم} به مرور پژوهش‌ها و کارهای مرتبط اختصاص دارد.
در این فصل پژوهش‌های مرتبط با تشخیص گفتار فارسی،
مدل‌های زبانی مورد استفاده برای زبان فارسی،
سامانه‌های تبدیل متن به گفتار فارسی و سامانه‌های مکالمه صوتی
مورد بررسی و مقایسه قرار می‌گیرند.

\textbf{فصل چهارم} طراحی و معماری سامانه پیشنهادی را تشریح
می‌کند.
در این فصل جریان پردازش از دریافت صوت کاربر تا تولید پاسخ صوتی،
نحوه ارتباط مؤلفه‌های
\lr{ASR}، \lr{LLM} و \lr{TTS}،
طراحی ماژولار سامانه، مدیریت بافت مکالمه و سایر پردازش‌های
تکمیلی توضیح داده می‌شوند.

\textbf{فصل پنجم} به جزئیات پیاده‌سازی اختصاص دارد.
در این فصل ابزارها و کتابخانه‌های مورد استفاده،
ساختار نرم‌افزاری پروژه، نحوه بارگذاری و اجرای مدل‌ها،
محیط اجرا، سخت‌افزار، رابط خط فرمان و ملاحظات مربوط به
بازتولیدپذیری تشریح می‌شوند.

\textbf{فصل ششم} روش آزمایش و نتایج تجربی پروژه را ارائه
می‌کند.
ابتدا ارزیابی مستقل مدل‌های تشخیص گفتار، مدل‌های زبانی و
تبدیل متن به گفتار گزارش شده و سپس عملکرد سامانه یکپارچه
از نظر تأخیر، کیفیت و مصرف منابع مورد بررسی قرار می‌گیرد.

\textbf{فصل هفتم} شامل بحث و تحلیل نتایج است.
در این فصل نقاط قوت و ضعف مدل‌های بررسی‌شده،
دلیل انتخاب مدل‌های نهایی، اثر خطاهای هر مرحله بر مراحل بعدی،
محدودیت‌های آزمایش‌ها و چالش‌های اجرای محلی سامانه تحلیل
خواهند شد.

در نهایت،
\textbf{فصل هشتم}
به نتیجه‌گیری، مرور دستاوردهای اصلی پروژه، محدودیت‌های باقی‌مانده
و پیشنهادهایی برای توسعه و پژوهش‌های آینده اختصاص دارد.
